% chapter1-body.tex — pure content, no preamble.
% This file is \input{} by both main.tex (book) and chapter1.tex (article).




\section{Problem Comprehend}
Consider a one-dimensional line with a wall located at
\[
    x=a>0.
\]
We launch a right-going wave from the left. Its complex representation is
\[
    U(x,t)=e^{ikx}e^{-i\omega t}
    =e^{i(kx-\omega t)}.
\]

Here, $\omega$ is the angular frequency, with unit $\mathrm{s}^{-1}$,
and $k$ is the wavenumber, with unit $\mathrm{m}^{-1}$.
The spatial phase is
\[
    \phi(x)=kx.
\]

For a uniform wave, one complete wavelength $\lambda$ corresponds to
a phase change of $2\pi$. Hence
\[
    k=\frac{2\pi}{\lambda}.
\]

This definition uses a complete wavelength. However, if the phase varies
spatially, a complete wavelength may not be locally available for
measurement. Instead, we consider the phase change over an infinitesimal
spatial interval:
\[
    k(x)
    =\lim_{\Delta x\to0}
    \frac{\phi(x+\Delta x)-\phi(x)}{\Delta x}
    =\frac{d\phi}{dx}.
\]

Thus, the local wavenumber measures the local phase accumulation per
unit distance.

The same idea can also be interpreted through time.
Dividing the local phase change and the propagation distance
by the same time interval gives, intuitively,
\[
    k
    =\frac{d\phi}{dx}
    \sim
    \frac{d\phi/dt}{dx/dt}
    =\frac{\omega}{c}.
\]
Thus, the wavenumber can also be understood as the temporal phase
accumulation rate divided by the propagation speed.

% TODO (Work 2): give a rigorous derivation using the total derivative
% along a constant-phase trajectory.

In higher dimensions, this idea becomes the local wavevector
\[
    \mathbf{k}(x)=\nabla\phi(x).
\]


\subsection{Propagation Directions in One Dimension}
In one dimension, the unit sphere consists of only two points,
\[
    S^0=\{-1,+1\},
\]
corresponding to the two possible propagation directions.

If
\[
    D^1=[-1,1],
\]
then
\[
    \partial D^1=S^0=\{-1,+1\}.
\]

For a direction $d\in S^0$, a plane-wave solution can be written as
\[
    u(x;d)=e^{ikxd}.
\]
Therefore,
\[
    d=+1 \quad\Longrightarrow\quad e^{ikx},
\]
and
\[
    d=-1 \quad\Longrightarrow\quad e^{-ikx}.
\]

% TODO: explain how the discrete directions in $S^0$
% generalize to $d\in S^{n-1}$ in higher dimensions.




\section{TMP}


\subsection{MVHR}
\[
    T^i{}_{ab}=
    \left[\begin{array}{c}
            \left[\begin{array}{ccc}
                          \left[\begin{array}{NNNN}
                                  1 &
                                  2 &
                                  3 &
                                  4
                              \end{array}\right]_{b}
                           &
                          \left[\begin{array}{NNNN}
                                  5 &
                                  6 &
                                  7 &
                                  8
                              \end{array}\right]_{b}
                           &
                          \left[\begin{array}{NNNN}
                                  9  &
                                  10 &
                                  11 &
                                  12
                              \end{array}\right]_{b}
                      \end{array}\right]_{a}
            \\
            \left[\begin{array}{ccc}
                          \left[\begin{array}{NNNN}
                                  13 &
                                  14 &
                                  15 &
                                  16
                              \end{array}\right]_{b}
                           &
                          \left[\begin{array}{NNNN}
                                  17 &
                                  18 &
                                  19 &
                                  20
                              \end{array}\right]_{b}
                           &
                          \left[\begin{array}{NNNN}
                                  21 &
                                  22 &
                                  23 &
                                  24
                              \end{array}\right]_{b}
                      \end{array}\right]_{a}
        \end{array}\right]_{i}
\]
\[
    T_i{}^a{}_b=
    \left[\begin{array}{cc}
            \left[\begin{array}{ccc}
                          \left[\begin{array}{NNNN}
                                  1 &
                                  2 &
                                  3 &
                                  4
                              \end{array}\right]_{b}
                          \\
                          \left[\begin{array}{NNNN}
                                  5 &
                                  6 &
                                  7 &
                                  8
                              \end{array}\right]_{b}
                          \\
                          \left[\begin{array}{NNNN}
                                  9  &
                                  10 &
                                  11 &
                                  12
                              \end{array}\right]_{b}
                      \end{array}\right]_{a}
             &
            \left[\begin{array}{ccc}
                          \left[\begin{array}{NNNN}
                                  13 &
                                  14 &
                                  15 &
                                  16
                              \end{array}\right]_{b}
                          \\
                          \left[\begin{array}{NNNN}
                                  17 &
                                  18 &
                                  19 &
                                  20
                              \end{array}\right]_{b}
                          \\
                          \left[\begin{array}{NNNN}
                                  21 &
                                  22 &
                                  23 &
                                  24
                              \end{array}\right]_{b}
                      \end{array}\right]_{a}
        \end{array}\right]_{i}
\]
\[
    s=\left[\begin{array}{c}7\end{array}\right]_{0}
\]
\[
    A^i=\left[\begin{array}{c}
            7
            \\
            3
        \end{array}\right]_{i}
\]
\[
    A_i=\left[\begin{array}{cc}7&3\end{array}\right]_{i}
\]
\[
    B^i{}_a=\left[\begin{array}{cc}
            \left[\begin{array}{lll}1&2&3\end{array}\right]_{a}
            \\
            \left[\begin{array}{lll}4&5&6\end{array}\right]_{a}
        \end{array}\right]_{i}
\]
\[
    B_i{}^a=\left[\begin{array}{cc}
            \left[\begin{array}{lll}
                          1
                          \\
                          2
                          \\
                          3\end{array}\right]_{a}
             &
            \left[\begin{array}{lll}
                          1
                          \\
                          2
                          \\
                          3\end{array}\right]_{a}
        \end{array}\right]_{i}
\]
\[
    B^{ia}=\left[\begin{array}{cc}
            \left[\begin{array}{lll}
                          1
                          \\
                          2
                          \\
                          3\end{array}\right]_{a}
            \\
            \left[\begin{array}{lll}
                          1
                          \\
                          2
                          \\
                          3\end{array}\right]_{a}
        \end{array}\right]_{i}
\]
\[
    B_{ia}=\left[\begin{array}{cc}
            \left[\begin{array}{lll}
                          1 &
                          2 &
                          3
                      \end{array}\right]_{a}
             &
            \left[\begin{array}{lll}
                          1 &
                          2 &
                          3
                      \end{array}\right]_{a}
        \end{array}\right]_{i}
\]
TMP




\section{Reference}
\bibliographystyle{plain}  % or another style like unsrt, IEEEtran, etc.
\bibliography{refer}  % references.bib is the file name


% \subsection{Definition of EMH}
% \begin{definition}[EMH means Symmetries]
%     \label{def:mydefinition}
%     We say a market is completely EMH if and
%     only if for every condition $A$, the random
%     variable of the logarithmic growth rate $W$
%     agrees that $\mathbb{E}_{W\mid A}(W)\doteq 0$.
% \end{definition}


% \subsection{Prior Work}
% Discussion of what others have done.




% \section{Methods}
% How you approached the problem.




% \section{Results}
% What you found.




% \section{Conclusion}
% Your closing thoughts.




